% Part: methods
% Chapter: induction
% Section: strong-induction

\documentclass[../../../include/open-logic-section]{subfiles}

\begin{document}

\olfileid[tr]{mth}{ind}{str}

\olsection{Kuvvetli Tümevarım}

Yukarıda ele alınan tümevarım ilkesinde $P(0)$'ı ve ayrıca $P(k)$ ise
$P(k+1)$ olduğunu kanıtlarız. İkinci kısımda $P(k)$'nin doğru olduğunu
varsayıp bu varsayımı $P(k+1)$'i kanıtlamak için kullanırız. Elbette
buna denk olarak $P(k-1)$'i varsayıp onu $P(k)$'yi kanıtlamak için
kullanabilirdik---önemli olan, herhangi bir sayıdan onun ardılına
çıkarımı gerçekleştirebilmemiz; söz konusu iddiayı, önceki sayı için
geçerli olduğu varsayımı altında herhangi bir sayı için
kanıtlayabilmemizdir.

Tümevarım ilkesinin bir çeşidinde, iddianın yalnızca $k$'nin önceki
sayısı $k-1$ için geçerli olduğunu değil,~$k$'den küçük bütün sayılar
için geçerli olduğunu varsayar ve bu varsayımı iddiayı~$k$ için
belirlemekte kullanırız. Bu da bize bütün~$n \in \Nat$ için $P(n)$
iddiasını verir. Çünkü $P(0)$'ı belirledikten sonra, $P$'nin~$1$'den
küçük bütün sayılar için geçerli olduğunu da belirlemiş oluruz. Ayrıca
bütün $l<k$ için $P(l)$ ise $P(k)$ olduğunu biliyorsak, bunu özellikle
$k=1$ için biliriz. Böylece $P(1)$ sonucuna varabiliriz. Bununla $P(0)$
ve $P(1)$'i, yani bütün $l<2$ için $P(l)$'yi kanıtlamış oluruz; bütün
$l<2$ için $P(l)$ ise $P(2)$ şeklindeki koşullu da elimizde olduğundan
$P(2)$ sonucuna varabiliriz ve bu böyle sürer.

Aslında ``bütün $k$'ler için, bütün $l<k$ için $P(l)$ ise $P(k)$''
genel koşullusunu belirleyebilirsek, artık $P(0)$'ı ayrıca belirlememiz
gerekmez; çünkü ondan çıkar. ``Bütün $l<k$ için $P(l)$'' gibi genel bir
iddianın, hiçbir $l<k$ yoksa doğru olduğunu hatırlayın. Bu, boş
nicelemenin bir örneğidir: hiç $A$ yoksa ``bütün $A$'lar $B$'dir''
doğrudur; hiçbir $x$,~$!A(x)$'i sağlamıyorsa
$\lforall[x][(!A(x) \lif !B(x))]$ doğrudur. Burada biçimselleştirilmiş
ifade ``$\lforall[l][(l < k \lif P(l))]$'' olur---ve hiç $l < k$ yoksa
bu doğrudur. $k=0$ olduğunda da durum tam olarak budur: hiçbir $l<0$
yoktur; dolayısıyla $P$ ne olursa olsun ``bütün $l<0$ için $P(l)$''
doğrudur. Böylece ``bütün $l<k$ için $P(l)$ ise $P(k)$''nin bir kanıtı
$P(0)$'ı kendiliğinden belirler.

Bu çeşit, iddiayı~$k$ için belirlemek yalnızca iddianın $k-1$ için
geçerli olmasına dayandırılamıyor, fakat onun bir ya da daha fazla
$l<k$ için doğru olduğu varsayımını gerektiriyorsa yararlıdır.

\end{document}
